%%
% BIThesis 本科毕业设计论文模板 —— 使用 XeLaTeX 编译 The BIThesis Template for Undergraduate Thesis
% This file has no copyright assigned and is placed in the Public Domain.
%%

\chapter{系统设计}

\section{总体架构设计}

本系统采用移动客户端、服务端、管理后台端三部分协同工作的总体架构。其中，移动客户端作为学习活动的主要承载体，负责课程展示、练习交互、挑战参与、社交反馈与个人成长呈现；服务端负责统一提供认证、课程内容、提交记录、奖励发放、社交关系和进度管理等业务能力；管理后台端负责课程、题目、挑战、用户与运营信息维护。三部分之间以统一的 OpenAPI 契约作为接口约束基础，从而保证数据结构、权限边界与交互流程的一致性。

为便于说明三端在系统中的职责分工，本文将其核心定位归纳如表\ref{tab:system-roles}所示。可以看出，移动客户端承担学习活动入口与交互反馈职能，服务端承担统一业务支撑与数据处理职能，管理后台端承担内容维护与运营配置职能，三者共同构成完整的学习闭环。

为了使三端之间的关系更加直观，论文终稿建议补充一张总体架构图，展示移动客户端、服务端、管理后台端以及 OpenAPI 契约文档之间的连接关系和主要数据流向。

\placeholderfigure{系统总体架构图占位}{fig:system-architecture}{待替换图片：系统总体架构图\\建议内容：移动客户端、服务端、管理后台端、OpenAPI 契约文档及其交互关系}

\begin{table}[htbp]
  \centering
  \caption{系统三端职责划分}
  \label{tab:system-roles}
  \begin{tabular}{p{3cm}p{4cm}p{6cm}}
    \toprule
    端类别 & 核心定位 & 主要职责 \\
    \midrule
    移动客户端 & 学习活动主载体 & 课程学习、章节阅读、编码练习、挑战参与、社交互动、成长反馈 \\
    服务端 & 统一业务支撑中心 & 认证鉴权、课程内容提供、提交记录、奖励结算、进度管理、社交数据处理 \\
    管理后台端 & 内容维护与运营支撑 & 课程管理、题目管理、挑战配置、用户管理、公告与运营配置 \\
    \bottomrule
  \end{tabular}
\end{table}

该架构的核心设计思想不是三端并列展开，而是围绕移动客户端学习闭环进行支撑式组织。也就是说，服务端与管理后台端的存在目的是保证移动客户端学习过程能够持续、稳定、可配置地运行，而不是与移动客户端分别形成独立的论文主体。在这种架构思路下，服务端被设计为统一的数据与业务处理中心，避免移动客户端直接访问数据库或执行复杂业务逻辑，从而保证数据安全性和业务规则的一致性。管理后台端则作为运营人员的操作界面，其所有操作最终都通过服务端接口生效，这种间接访问方式使系统能够在管理后台端和移动客户端之间建立统一的业务规则执行层。

从部署角度看，服务端作为独立服务运行，管理后台端作为静态前端应用部署，移动客户端则作为原生应用分发到用户设备。这种分离部署方式使各端可以独立迭代和发布，减少了版本耦合带来的协调成本。同时，服务端统一对外提供接口，使移动客户端和管理后台端共享同一套业务逻辑和数据模型，降低了多端维护的复杂度。

\section{功能结构设计}

根据移动客户端路由与页面组织，系统移动客户端功能结构可划分为启动与认证、课程学习、编码练习、闯关挑战、每日挑战、社交互动、个人中心七个部分。启动与认证负责完成用户状态识别与基础登录流程；课程学习负责承载课程目录、课程详情与章节内容；编码练习负责连接题目内容、代码提交与结果反馈；闯关挑战与每日挑战用于组织具有目标性的任务链；社交互动负责呈现好友、动态与排行榜；个人中心负责展示用户资料、学习统计和奖励信息。

在启动与认证模块中，系统需要处理应用冷启动时的初始化流程，包括本地状态检查、登录态恢复判断和引导页面展示。如果本地存在可用的认证信息，系统可直接恢复到主流程界面；如果认证信息失效或不存在，则引导用户进入登录或注册流程。这种设计使已登录用户能够较快恢复学习状态，同时保证未登录用户能够顺利进入系统。

在课程学习模块中，功能结构进一步细分为课程列表、课程详情和章节阅读三个层次。课程列表提供全部课程的概览信息和筛选能力，课程详情展示课程介绍、章节目录和相关练习，章节阅读则承载具体的知识内容。这种三层结构使学习者能够从宏观到微观逐步深入，同时保持每一层的信息密度在移动设备可接受范围内。

在编码练习模块中，功能结构围绕题目展示、代码编辑和结果反馈三个核心环节展开。题目展示需要清晰呈现题目要求、示例输入输出和注意事项；代码编辑需要提供基础的输入环境和语法辅助；结果反馈需要及时返回编译结果、运行输出和测试用例通过情况。这三个环节构成了练习功能的完整闭环，任何一环的体验缺陷都会影响学习者的练习积极性。

与之对应，管理后台端支撑结构主要包含课程管理、挑战管理、题目管理、用户管理、内容审核和公告配置等能力。管理后台端并不直接面向学习活动本身，而是通过维护内容与配置规则，为移动客户端持续提供可消费的数据与管理能力。在课程管理中，运营人员可以创建和编辑课程信息、组织章节结构、关联练习题目和控制课程上下线状态。在挑战管理中，运营人员可以配置关卡顺序、设置通关条件、调整奖励规则和管理挑战可用性。在用户管理中，运营人员可以查看用户基础信息、学习统计和行为记录，为运营决策提供数据支持。

在功能结构表达上，建议补充一张移动客户端功能结构图，以更清楚地对应课程学习、编码练习、挑战任务、社交互动与个人中心等模块之间的组织关系。

\placeholderfigure{移动客户端功能结构图占位}{fig:mobile-function-structure}{待替换图片：移动客户端功能结构图\\建议内容：启动与认证、课程学习、编码练习、闯关挑战、每日挑战、社交互动、个人中心}

\section{核心业务流程设计}

在核心学习流程上，学习者首先通过认证模块进入系统，随后浏览课程列表并进入具体课程与章节，在完成知识阅读后进入编码练习或挑战任务。练习与挑战的提交结果会驱动学习进度、经验值、奖励和统计数据更新，并进一步影响排行榜、个人中心与社交动态展示。这样的设计使得知识学习、行为记录和激励反馈形成闭环，从而增强学习过程的连续性。

具体而言，学习流程的起点是认证模块。学习者输入账号密码完成登录，或使用注册功能创建新账号。登录成功后，服务端返回认证 Token，移动客户端将其持久化到本地存储，并在后续请求中自动携带。服务端同时提供 Token 刷新接口，客户端当前主要在登录失效时进行统一处理并引导重新登录，从而维持基础会话管理能力。

进入主界面后，学习者首先接触的是课程列表。课程列表展示标题、简介、进度等基础信息，并支持搜索与筛选。学习者选择感兴趣的课程后进入课程详情页，查看课程介绍和章节目录。章节按照顺序排列，已完成的章节标记为完成状态，未完成的章节中第一个未解锁的章节标记为当前学习进度。这种视觉反馈使学习者能够直观了解自己的学习位置和剩余内容。

在章节阅读页面，学习者浏览知识内容，内容以图文形式组织，适应移动设备的阅读习惯。阅读过程中，系统结合章节完成状态与课程进度记录学习过程。完成章节阅读后，学习者可以选择进入关联的练习页面，通过编码实践巩固所学知识。练习页面展示题目要求和代码编辑区，学习者编写代码后提交，服务端对代码进行评测并返回结果反馈。根据反馈，学习者可以修改代码重新提交，逐步完成练习目标。

在每日挑战流程中，系统按照每日任务入口向学习者提供限时或周期性任务，学习者完成提交后更新打卡记录、连续学习状态与额外奖励。该流程与普通挑战相比，更强调节奏控制与日常激励，从而提高用户的持续访问频率。每日挑战的设计考虑了学习者的时间安排习惯，任务难度适中，完成时间可控，使学习者能够在碎片时间内完成并获得成就感。连续打卡机制则通过累积奖励激励学习者保持日常学习习惯，但系统设计时也考虑了断签后的恢复机制，避免因偶然中断造成过度挫败感。

为了帮助读者理解学习活动如何从登录进入课程、练习、挑战再回到成长反馈，建议在此处补充一张核心学习流程图。

\placeholderfigure{核心学习流程图占位}{fig:learning-flow}{待替换图片：核心学习流程图\\建议内容：登录/注册、课程学习、章节阅读、练习提交、挑战参与、奖励反馈、排行榜/个人中心}

\section{接口与数据设计}

为了保证多端协同一致性，系统将 OpenAPI 文档作为共享接口真源，对路径、方法、请求体、响应结构、受众标记与权限边界进行统一定义。服务端实现依据该接口文档组织路由与处理逻辑，移动客户端通过统一的请求服务访问接口并完成 Token 注入、异常处理与结果解包，管理后台端则通过统一 API 封装消费管理能力。

在接口设计原则上，系统遵循 RESTful 风格，使用 HTTP 方法表达操作语义，使用路径参数和查询参数传递定位信息，使用请求体传递复杂数据。响应结构采用统一的包络格式，包含状态码、消息文本和数据载荷三个部分，使调用方能够以一致的方式处理成功和失败情况。对于列表类接口，系统支持分页参数，使移动客户端能够按需加载数据，避免一次性传输大量内容。

在数据结构上，系统围绕课程、章节、练习、挑战、每日挑战、用户资料、好友关系、成长奖励与学习进度等对象进行设计。其中，学习者侧的数据对象需要兼顾展示需求与交互需求，而服务端数据结构则更强调持久化与状态更新逻辑。通过契约文档、状态机说明与字段词典的配合，可以在设计阶段约束字段语义与状态流转，降低实现阶段出现理解偏差的可能性。

课程对象包含课程标识、标题、简介、封面图、分类、难度等级、上下线状态和创建时间等字段。章节对象包含章节标识、所属课程、标题、内容、顺序号和关联练习标识等字段。练习对象包含练习标识、题目描述、示例输入输出、测试用例、难度等级和关联章节等字段。挑战对象包含挑战标识、标题、描述、关卡列表、奖励规则和可用性状态等字段。用户资料对象包含用户标识、昵称、头像、等级、经验值、注册时间和最近活动时间等字段。

在状态设计方面，系统对关键实体定义了明确的状态集合和流转规则。例如，课程具有草稿、已上线和已下线三种状态，状态转换由管理后台端操作触发。挑战参与记录具有未开始、进行中、已完成和已放弃四种状态，状态转换由学习者的操作触发。这些状态定义和流转规则在契约文档中以状态机形式描述，为多端实现提供了一致的理解基础。

在接口与数据设计部分，建议补充一张接口协同或关键数据流图，用于呈现客户端请求、服务端处理以及管理后台端配置变更之间的关系。

\placeholderfigure{接口协同关系图占位}{fig:api-collaboration}{待替换图片：接口协同关系图\\建议内容：移动客户端请求、服务端处理、管理后台端配置、数据库或契约文档关系}

\section{服务端与管理后台端协同设计}

服务端路由按照学习者侧与管理侧受众进行组织，其中学习者相关接口主要服务于课程、资料、进度、挑战、奖励、社交与提交流程；管理侧接口则服务于课程、挑战、题目、用户、反馈与配置维护。这样的划分方式能够使系统在权限控制与职责边界上更加清晰，也有利于后续联调测试中的问题定位。

学习者侧接口采用统一的认证中间件进行身份校验，从请求头中提取 Token 并解析用户身份，将身份信息附加到请求上下文中供后续处理使用。这种设计使业务处理逻辑无需关心认证细节，只需从上下文中获取已验证的用户身份即可。管理侧接口除了认证校验外，还增加了角色校验，确保只有具备管理权限的用户才能访问管理功能。

在路由组织上，学习者侧接口以学习者视角组织资源路径，例如课程列表、我的进度、我的挑战等；管理侧接口则以管理视角组织资源路径，例如课程管理、用户管理、统计报表等。这种路径命名方式使接口的用途和受众一目了然，减少了调用方的理解成本。

管理后台端通过认证守卫、统一 API 访问层和页面级管理入口完成支撑工作，主要用于内容维护与运营配置。管理后台端主要检查登录状态，未登录用户会被重定向到登录页面；更具体的管理权限约束主要由服务端接口控制。统一 API 访问层封装了所有与服务端的通信逻辑，包括请求发送、错误处理和响应解析，使页面组件只需关注业务逻辑而无需处理网络细节。

由于本文以移动客户端学习应用为主体，因此管理后台端设计在论文中只说明其对内容供给与管理闭环的支撑意义，而不对各个管理页面展开细粒度实现分析。但需要强调的是，管理后台端的完备性直接影响移动客户端的内容质量和运营效率。如果管理后台端缺乏灵活的内容编辑能力，运营人员将难以快速响应内容更新需求；如果管理后台端缺乏有效的数据查看能力，运营决策将缺乏数据支撑。因此，在系统设计时，管理后台端的功能范围与易用性被纳入整体考虑，确保其能够切实支撑移动客户端的持续运营。

在数据流向上，管理后台端的操作通过服务端接口持久化到数据库，移动客户端通过服务端接口读取最新数据。这种单向数据流设计使数据变更的源头唯一且可追溯，避免了多端同时修改数据可能带来的冲突和不一致。在缓存策略上，移动客户端对课程列表等相对静态的数据进行本地缓存，减少重复请求；对进度、排行榜等动态数据则优先从服务端获取，保证信息的时效性。

在安全设计方面，系统从接口访问控制、数据传输保护和输入验证三个层面保障安全性。接口访问控制通过 Token 认证和角色校验实现，确保只有合法用户能够访问其权限范围内的资源。数据传输保护在部署层面可结合 HTTPS 等方式进行实现，以降低敏感信息在传输过程中被截获或篡改的风险。服务端对关键输入进行了基础类型校验和必要的业务校验，减少异常输入导致的数据错误。虽然本文不涉及专业的安全审计和渗透测试，但这些基础安全措施为系统的稳定运行提供了必要的保障。

在错误处理设计方面，系统建立了分层的错误处理机制。服务端对业务错误、认证错误和系统错误进行分类，分别返回对应的 HTTP 状态码和错误信息。移动客户端根据错误类型采取不同的处理策略：对于网络错误，提示用户检查网络连接并提供重试选项；对于认证错误，引导用户重新登录；对于业务错误，展示服务端返回的具体错误信息。这种分层处理机制使错误信息对用户更加友好，同时也便于开发人员进行问题排查和定位。

在日志与监控设计方面，服务端记录了基础请求日志信息，包括请求路径、响应状态、处理时间和异常详情等。这些日志信息在开发和测试阶段用于问题定位，在运行阶段也可作为观察系统状态的基础。移动客户端当前主要保留基础错误输出与状态提示，用于辅助调试与用户反馈。虽然本文未实现完整的分布式追踪和实时监控体系，但基础的日志记录为系统的可观测性提供了初步支持。

在数据库设计方面，系统采用关系型数据库进行数据持久化，通过合理的表结构设计和索引策略保证基础查询性能。核心数据表包括用户表、课程表、章节表、练习表、挑战表、提交记录表、进度表、好友关系表和奖励记录表等。表之间通过外键关联建立数据一致性约束，同时避免过度规范化导致的查询复杂度上升。在数据库访问层，系统使用连接池管理数据库连接，通过预编译语句防止 SQL 注入攻击，这些基础实践为数据安全和访问性能提供了保障。
